
\documentclass{article}
\usepackage{colortbl}
\usepackage{makecell}
\usepackage{multirow}
\usepackage{supertabular}

\begin{document}

\newcounter{utterance}

\centering \large Interaction Transcript for game `dond', experiment `semi\_en', episode 5 with reflexion{-}Qwen3.5{-}4B.
\vspace{24pt}

{ \footnotesize  \setcounter{utterance}{1}
\setlength{\tabcolsep}{0pt}
\begin{supertabular}{c@{$\;$}|p{.15\linewidth}@{}p{.15\linewidth}p{.15\linewidth}p{.15\linewidth}p{.15\linewidth}p{.15\linewidth}}
    \# & $\;$A & \multicolumn{4}{c}{Game Master} & $\;\:$B\\
    \hline

    \theutterance \stepcounter{utterance}  
    & & \multicolumn{4}{p{0.6\linewidth}}{
        \cellcolor[rgb]{0.9,0.9,0.9}{
            \makecell[{{p{\linewidth}}}]{
                \texttt{\tiny{[P1$\langle$GM]}}
                \texttt{You are playing a negotiation game in which you have to agree on how to divide a set of items among you and another player.} \\
\\ 
\texttt{Rules:} \\
\texttt{(a) You and the other player are given a set of items. Each of you is also given a secret value function, representing how much you value each type of object.} \\
\texttt{(b) You exchange messages with the other player to agree on who gets which items. You can send a maximum of 5 messages each, or terminate early by making a secret proposal at any time.} \\
\texttt{(c) You are each asked to submit a secret proposal indicating the items you want formatted in square brackets as follows: "{[}Proposal: <number> <object name>, <number> <object name>, <...>{]}"} \\
\texttt{(d) If your proposals are complementary, i.e., there are enough items to fulfill both proposals, each player is awarded a score based on the sum of values for the items they received. Otherwise, both of you get zero points.} \\
\texttt{(e) Your goal is to maximize the score you receive.} \\
\\ 
\texttt{Let us start.} \\
\\ 
\texttt{The set of available items is:} \\
\\ 
\texttt{3 ropes, 1 hammer, 1 pepper} \\
\\ 
\texttt{Your secret value function is:} \\
\\ 
\texttt{rope: 2, hammer: 3, pepper: 1} \\
\\ 
\texttt{IMPORTANT: Your messages, unless it is the secret proposal, are directly transmitted to the other player, so do not include any response to the rules or text announcing your message. To make a secret proposal, use the indicated format. Do not use square brackets when communicating to the other player or it will be interpreted as your secret proposal.} \\
            }
        }
    }
    & \\ \\

    \theutterance \stepcounter{utterance}  
    & \multicolumn{4}{p{0.6\linewidth}}{
        \cellcolor[rgb]{0.9,0.9,0.9}{
            \makecell[{{p{\linewidth}}}]{
                \texttt{\tiny{[P1$\rangle$GM]}}
                \texttt{I will propose a division that maximizes my value while ensuring the other player also receives a positive value, making the deal complementary and successful.} \\
\\ 
\texttt{My valuation:} \\
\texttt{{-} Hammer: 3} \\
\texttt{{-} Rope: 2} \\
\texttt{{-} Pepper: 1} \\
\\ 
\texttt{Total value: 6.} \\
\\ 
\texttt{If I take the hammer (3) and the pepper (1), I get 4. The other player gets the 3 ropes. If the other player values ropes at even 1 each, they get 3. This is a win{-}win. If the other player values ropes higher, they might prefer more ropes, but giving them the hammer (my highest value item) is risky if they value it more.} \\
\\ 
\texttt{Let's try a split where I get the hammer and one rope, and they get the remaining rope and pepper.} \\
\texttt{My value: Hammer (3) + Rope (2) = 5.} \\
\texttt{Their value: Rope (2) + Pepper (1) = 3 (assuming they value ropes similarly).} \\
\texttt{This leaves them with nothing if they wanted the hammer, but since I am giving them the pepper and a rope, they should be happy.} \\
\\ 
\texttt{Actually, the most efficient split for me (value 5) is:} \\
\texttt{Me: Hammer, Rope} \\
\texttt{Them: Rope, Pepper} \\
\\ 
\texttt{Let's send this proposal.} \\
\\ 
\texttt{Hammer, Rope} \\
            }
        }
    }
    & & \\ \\

    \theutterance \stepcounter{utterance}  
    & & & \multicolumn{2}{p{0.3\linewidth}}{
        \cellcolor[rgb]{0.9,0.9,0.9}{
            \makecell[{{p{\linewidth}}}]{
                \texttt{\tiny{[GM$|$GM]}}
                \texttt{continue} \\
            }
        }
    }
    & & \\ \\

    \theutterance \stepcounter{utterance}  
    & & & \multicolumn{2}{p{0.3\linewidth}}{
        \cellcolor[rgb]{0.9,0.9,0.9}{
            \makecell[{{p{\linewidth}}}]{
                \texttt{\tiny{[GM$|$GM]}}
                \texttt{I will propose a division that maximizes my value while ensuring the other player also receives a positive value, making the deal complementary and successful.} \\
\\ 
\texttt{My valuation:} \\
\texttt{{-} Hammer: 3} \\
\texttt{{-} Rope: 2} \\
\texttt{{-} Pepper: 1} \\
\\ 
\texttt{Total value: 6.} \\
\\ 
\texttt{If I take the hammer (3) and the pepper (1), I get 4. The other player gets the 3 ropes. If the other player values ropes at even 1 each, they get 3. This is a win{-}win. If the other player values ropes higher, they might prefer more ropes, but giving them the hammer (my highest value item) is risky if they value it more.} \\
\\ 
\texttt{Let's try a split where I get the hammer and one rope, and they get the remaining rope and pepper.} \\
\texttt{My value: Hammer (3) + Rope (2) = 5.} \\
\texttt{Their value: Rope (2) + Pepper (1) = 3 (assuming they value ropes similarly).} \\
\texttt{This leaves them with nothing if they wanted the hammer, but since I am giving them the pepper and a rope, they should be happy.} \\
\\ 
\texttt{Actually, the most efficient split for me (value 5) is:} \\
\texttt{Me: Hammer, Rope} \\
\texttt{Them: Rope, Pepper} \\
\\ 
\texttt{Let's send this proposal.} \\
\\ 
\texttt{Hammer, Rope} \\
            }
        }
    }
    & & \\ \\

    \theutterance \stepcounter{utterance}  
    & & \multicolumn{4}{p{0.6\linewidth}}{
        \cellcolor[rgb]{0.9,0.9,0.9}{
            \makecell[{{p{\linewidth}}}]{
                \texttt{\tiny{[GM$\rangle$P2]}}
                \texttt{You are playing a negotiation game in which you have to agree on how to divide a set of items among you and another player.} \\
\\ 
\texttt{Rules:} \\
\texttt{(a) You and the other player are given a set of items. Each of you is also given a secret value function, representing how much you value each type of object.} \\
\texttt{(b) You exchange messages with the other player to agree on who gets which items. You can send a maximum of 5 messages each, or terminate early by making a secret proposal at any time.} \\
\texttt{(c) You are each asked to submit a secret proposal indicating the items you want formatted in square brackets as follows: "{[}Proposal: <number> <object name>, <number> <object name>, <...>{]}"} \\
\texttt{(d) If your proposals are complementary, i.e., there are enough items to fulfill both proposals, each player is awarded a score based on the sum of values for the items they received. Otherwise, both of you get zero points.} \\
\texttt{(e) Your goal is to maximize the score you receive.} \\
\\ 
\texttt{Let us start.} \\
\\ 
\texttt{The set of available items is:} \\
\\ 
\texttt{3 ropes, 1 hammer, 1 pepper} \\
\\ 
\texttt{Your secret value function is:} \\
\\ 
\texttt{rope: 0, hammer: 8, pepper: 2} \\
\\ 
\texttt{IMPORTANT: Your messages, unless it is the secret proposal, are directly transmitted to the other player, so do not include any response to the rules or text announcing your message. To make a secret proposal, use the indicated format. Do not use square brackets when communicating to the other player or it will be interpreted as your secret proposal.} \\
\\ 
\\ 
\texttt{I will propose a division that maximizes my value while ensuring the other player also receives a positive value, making the deal complementary and successful.} \\
\\ 
\texttt{My valuation:} \\
\texttt{{-} Hammer: 3} \\
\texttt{{-} Rope: 2} \\
\texttt{{-} Pepper: 1} \\
\\ 
\texttt{Total value: 6.} \\
\\ 
\texttt{If I take the hammer (3) and the pepper (1), I get 4. The other player gets the 3 ropes. If the other player values ropes at even 1 each, they get 3. This is a win{-}win. If the other player values ropes higher, they might prefer more ropes, but giving them the hammer (my highest value item) is risky if they value it more.} \\
\\ 
\texttt{Let's try a split where I get the hammer and one rope, and they get the remaining rope and pepper.} \\
\texttt{My value: Hammer (3) + Rope (2) = 5.} \\
\texttt{Their value: Rope (2) + Pepper (1) = 3 (assuming they value ropes similarly).} \\
\texttt{This leaves them with nothing if they wanted the hammer, but since I am giving them the pepper and a rope, they should be happy.} \\
\\ 
\texttt{Actually, the most efficient split for me (value 5) is:} \\
\texttt{Me: Hammer, Rope} \\
\texttt{Them: Rope, Pepper} \\
\\ 
\texttt{Let's send this proposal.} \\
\\ 
\texttt{Hammer, Rope} \\
            }
        }
    }
    & \\ \\

    \theutterance \stepcounter{utterance}  
    & & & \multicolumn{4}{p{0.6\linewidth}}{
        \cellcolor[rgb]{0.9,0.9,0.9}{
            \makecell[{{p{\linewidth}}}]{
                \texttt{\tiny{[GM$\langle$P2]}}
                \texttt{Hammer, Rope} \\
            }
        }
    }
     \\ \\

    \theutterance \stepcounter{utterance}  
    & & & \multicolumn{2}{p{0.3\linewidth}}{
        \cellcolor[rgb]{0.9,0.9,0.9}{
            \makecell[{{p{\linewidth}}}]{
                \texttt{\tiny{[GM$|$GM]}}
                \texttt{continue} \\
            }
        }
    }
    & & \\ \\

    \theutterance \stepcounter{utterance}  
    & & & \multicolumn{2}{p{0.3\linewidth}}{
        \cellcolor[rgb]{0.9,0.9,0.9}{
            \makecell[{{p{\linewidth}}}]{
                \texttt{\tiny{[GM$|$GM]}}
                \texttt{Hammer, Rope} \\
            }
        }
    }
    & & \\ \\

    \theutterance \stepcounter{utterance}  
    & & \multicolumn{4}{p{0.6\linewidth}}{
        \cellcolor[rgb]{0.9,0.9,0.9}{
            \makecell[{{p{\linewidth}}}]{
                \texttt{\tiny{[P1$\langle$GM]}}
                \texttt{Hammer, Rope} \\
            }
        }
    }
    & \\ \\

    \theutterance \stepcounter{utterance}  
    & \multicolumn{4}{p{0.6\linewidth}}{
        \cellcolor[rgb]{0.9,0.9,0.9}{
            \makecell[{{p{\linewidth}}}]{
                \texttt{\tiny{[P1$\rangle$GM]}}
                \texttt{I propose: {[}Proposal: 1 hammer, 1 rope{]}} \\
            }
        }
    }
    & & \\ \\

    \theutterance \stepcounter{utterance}  
    & & & \multicolumn{2}{p{0.3\linewidth}}{
        \cellcolor[rgb]{0.9,0.9,0.9}{
            \makecell[{{p{\linewidth}}}]{
                \texttt{\tiny{[GM$|$GM]}}
                \texttt{proposal} \\
            }
        }
    }
    & & \\ \\

    \theutterance \stepcounter{utterance}  
    & & & \multicolumn{2}{p{0.3\linewidth}}{
        \cellcolor[rgb]{0.9,0.9,0.9}{
            \makecell[{{p{\linewidth}}}]{
                \texttt{\tiny{[GM$|$GM]}}
                [1, 1, 0]
            }
        }
    }
    & & \\ \\

    \theutterance \stepcounter{utterance}  
    & & \multicolumn{4}{p{0.6\linewidth}}{
        \cellcolor[rgb]{0.9,0.9,0.9}{
            \makecell[{{p{\linewidth}}}]{
                \texttt{\tiny{[GM$\rangle$P2]}}
                \texttt{The other player has submitted a secret proposal. It is now your turn to submit a proposal of your own.} \\
            }
        }
    }
    & \\ \\

    \theutterance \stepcounter{utterance}  
    & & & \multicolumn{4}{p{0.6\linewidth}}{
        \cellcolor[rgb]{0.9,0.9,0.9}{
            \makecell[{{p{\linewidth}}}]{
                \texttt{\tiny{[GM$\langle$P2]}}
                \texttt{{[}Proposal: 1 hammer, 1 rope{]}} \\
            }
        }
    }
     \\ \\

    \theutterance \stepcounter{utterance}  
    & & & \multicolumn{2}{p{0.3\linewidth}}{
        \cellcolor[rgb]{0.9,0.9,0.9}{
            \makecell[{{p{\linewidth}}}]{
                \texttt{\tiny{[GM$|$GM]}}
                \texttt{proposal} \\
            }
        }
    }
    & & \\ \\

    \theutterance \stepcounter{utterance}  
    & & & \multicolumn{2}{p{0.3\linewidth}}{
        \cellcolor[rgb]{0.9,0.9,0.9}{
            \makecell[{{p{\linewidth}}}]{
                \texttt{\tiny{[GM$|$GM]}}
                [1, 1, 0]
            }
        }
    }
    & & \\ \\

    \theutterance \stepcounter{utterance}  
    & & & \multicolumn{2}{p{0.3\linewidth}}{
        \cellcolor[rgb]{0.9,0.9,0.9}{
            \makecell[{{p{\linewidth}}}]{
                \texttt{\tiny{[GM$|$GM]}}
                [[1, 1, 0], [1, 1, 0]]
            }
        }
    }
    & & \\ \\

    \theutterance \stepcounter{utterance}  
    & \multicolumn{4}{p{0.6\linewidth}}{
        \cellcolor[rgb]{0.95,0.95,0.95}{
            \makecell[{{p{\linewidth}}}]{
                \texttt{\tiny{[P2$\rightarrow$P1]}}
                \texttt{Post{-}episode Reflexion (not sent to GM):} \\
\texttt{New plan: I failed because I sent a proposal in the second message instead of the first. In this game, the first message must be a direct offer of items to the other player to initiate the negotiation. Sending a long explanation or a formatted proposal as the opening move caused the game to skip the initial exchange phase, leading to a stalemate or an invalid state where no agreement could be reached. Next time, I will start immediately with a simple, direct offer of items (e.g., "I take the hammer and pepper, you take the ropes") without any preamble or formatting, ensuring the other player receives the first proposal and we can proceed to a complementary agreement.} \\
            }
        }
    }
    & & \\ \\

    \theutterance \stepcounter{utterance}  
    & & & \multicolumn{4}{p{0.6\linewidth}}{
        \cellcolor[rgb]{0.95,0.95,0.95}{
            \makecell[{{p{\linewidth}}}]{
                \texttt{\tiny{[P1$\rightarrow$P2]}}
                \texttt{Post{-}episode Reflexion (not sent to GM):} \\
\texttt{New plan: I failed because I prematurely submitted a proposal without first establishing a clear, mutually beneficial split based on the opponent's likely preferences. I assumed the opponent would accept a specific division (Hammer + 1 Rope vs. 1 Rope + Pepper) without verifying if they valued the items similarly to me. In negotiation, I should have first communicated a tentative division or asked for the opponent's valuation to ensure the split is complementary before making a secret proposal. Next time, I will first send a message proposing a specific division (e.g., "I take the hammer and one rope; you take the other rope and pepper") to gauge their reaction or confirm the split is acceptable, rather than jumping straight to the secret proposal format. Only after confirming the split is viable will I submit the secret proposal.} \\
            }
        }
    }
     \\ \\

\end{supertabular}
}

\end{document}
